


The UniZ portal is your central hub for academics, leave requests, and your student profile. Your login credentials are your RGUKT roll number and a password set by the administration.

\subsection{Sign in}


\section{Instruction Step}

  
\section{Instruction Step}

    Navigate to the UniZ student portal URL provided by your institution.
  
  
\section{Instruction Step}

    - **Username** — your roll number (e.g., `O210008`) - **Password** — your
    current password
  
  
\section{Instruction Step}

    On success, you are redirected to your profile dashboard. The portal stores
    your session securely in the browser — you do not need to enter your
    credentials again until you log out or your session expires.
  
</Steps>

<Note>
  Your username is always your roll number in uppercase (e.g., `O210008`). It
  never changes.
</Note>

\subsubsection{How authentication works}

When you sign in, the server issues a JWT (JSON Web Token) that the portal stores in your browser. Every subsequent request automatically includes this token — you don't need to do anything extra in the portal.

If you are building on top of the API directly, include the token in every request:

\begin{lstlisting}[language=bash]
Authorization: Bearer <your-token>

\end{lstlisting}

**Login endpoint (for API use):**

\begin{lstlisting}[language=bash]
POST https://api.uniz.rguktong.in/api/v1/auth/login/student

\end{lstlisting}

\begin{lstlisting}[language=json]
{
  "username": "O210008",
  "password": "password123"
}

\end{lstlisting}

**Success response:**

\begin{lstlisting}[language=json]
{
  "success": true,
  "token": "eyJhbGciOiJIUz...",
  "role": "student",
  "username": "O210008"
}

\end{lstlisting}

---

\subsection{Reset your password}

Use this flow if you've forgotten your password. You'll need access to the email address registered with your account.


\section{Instruction Step}

  
\section{Instruction Step}

    Tap **Forgot Password** on the sign-in screen.
  
  
\section{Instruction Step}

    Submit your username (`O210008`). The system sends a 6-digit OTP to your registered email address.

    \begin{lstlisting}[language=json]
    POST /auth/otp/request
    { "username": "O210008" }
    
\end{lstlisting}

  
  
\section{Instruction Step}

    Check your inbox and enter the 6-digit code in the verification field.

    \begin{lstlisting}[language=json]
    POST /auth/otp/verify
    {
      "username": "O210008",
      "otp": "123456"
    }
    
\end{lstlisting}

  
  
\section{Instruction Step}

    Enter and confirm your new password. The reset requires a valid OTP and the same username.

    \begin{lstlisting}[language=json]
    POST /auth/password/reset
    {
      "username": "O210008",
      "resetToken": "123456",
      "newPassword": "newSecurePassword123"
    }
    
\end{lstlisting}

  
</Steps>

<Warning>
  OTP requests are rate-limited. If you request too many OTPs in a short period,
  you'll be temporarily blocked. Wait a few minutes before trying again.
</Warning>

---

\subsection{Change your password}

Once logged in, you can change your password from the **Reset Password** page in the portal. You'll need your current password.

A strong password must meet all three criteria:

- At least 8 characters
- Contains at least one number
- Contains at least one special character (e.g., `!@#$%`)

<Warning>
  After a successful password change, all active sessions are immediately logged
  out for security. You'll need to sign in again with your new password.
</Warning>

The portal evaluates your new password in real time and shows a strength indicator (Weak / Moderate / Strong). Weak and Moderate passwords are rejected.

---

\subsection{Install as a PWA}

UniZ is a Progressive Web App (PWA). You can install it on your phone or desktop so it opens like a native app — no app store required.

<Tabs>
  <Tab title="Android">
    
\section{Instruction Step}

      
\section{Instruction Step}

        Navigate to the UniZ portal URL.
      
      
\section{Instruction Step}

        Chrome shows an **Add to Home Screen** banner at the bottom, or you can
        tap the three-dot menu and select **Install app** / **Add to Home
        Screen**.
      
      
\section{Instruction Step}

        Tap **Install**. The app icon appears on your home screen and opens in
        standalone mode.
      
    </Steps>
  </Tab>
  <Tab title="iOS (Safari)">
    
\section{Instruction Step}

      
\section{Instruction Step}

        PWA installation on iOS only works in Safari.
      
      
\section{Instruction Step}

        Tap the share icon at the bottom of the screen (the box with an arrow
        pointing up).
      
      
\section{Instruction Step}

        Scroll down in the share sheet and tap **Add to Home Screen**, then tap
        **Add**.
      
    </Steps>
  </Tab>
  <Tab title="Desktop (Chrome/Edge)">
    
\section{Instruction Step}

      
\section{Instruction Step}

        Navigate to the UniZ portal URL in Chrome or Edge.
      
      
\section{Instruction Step}

        Look for an install icon in the address bar (a monitor with a download
        arrow), or open the browser menu and select **Install UniZ**.
      
      
\section{Instruction Step}

        Click **Install**. The app opens in its own window.
      
    </Steps>
  </Tab>
</Tabs>

<Tip>
  Installing the PWA gives you a faster experience and keeps your session active
  between visits without navigating to the URL each time.
</Tip>

---

\subsection{Common errors}


\begin{tcolorbox}[colback=gray!5,colframe=primary,title=Documentation Note]

  
\begin{tcolorbox}[colback=gray!5,colframe=primary,title=Documentation Note]

    Double-check that your roll number is entered in the correct format (e.g.,
    `O210008` — capital letter O, not zero). Passwords are case-sensitive.
  
\end{tcolorbox}

  
\begin{tcolorbox}[colback=gray!5,colframe=primary,title=Documentation Note]

    If your account has been suspended by the administration, you'll see a
    suspension message on sign-in. Contact your warden or SWO to resolve this.
  
\end{tcolorbox}

  
\begin{tcolorbox}[colback=gray!5,colframe=primary,title=Documentation Note]

    Check your spam folder. If the email still hasn't arrived after a few
    minutes, confirm your registered email with your hostel admin and try
    requesting the OTP again.
  
\end{tcolorbox}

  
\begin{tcolorbox}[colback=gray!5,colframe=primary,title=Documentation Note]

    JWT tokens expire after a set period. If you're redirected to the login page
    unexpectedly, simply sign in again.
  
\end{tcolorbox}

</CardGroup>


\section{Deep Technical Analysis}
The underlying system architecture for this module is built upon the \textbf{UniZ Microservices Framework}. 
This design ensures that each functional unit (Auth, Profile, Academics) remains isolated, preventing cascading failures across the university ecosystem.

\subsection{Operational Hardening}
Deployments are managed via K3s (Kubernetes) and containerized using high-performance Alpine Linux Docker images. 
Resource management is critical; for instance, the documentation service itself is provisioned with 2GB of RAM to ensure the responsive rendering of complex documentation graphs and state-management components.

\subsection{Enterprise Security Topology}
Every request entering the UniZ gateway is subjected to an aggressive JWT validation layer. 
Permissions are granular, mapping directly onto the RBAC (Role-Based Access Control) matrix defined in the core system specifications.

\subsection{Data Governance}
Prisma-driven data access ensures type safety and predictable query performance. 
We utilize PostgreSQL connection pooling to handle concurrent requests from the student portal and administrative dashboards simultaneously.
